\section{Effective Sample Size}
\label{sec:ess}

\subsection{Definition}

A Markov chain of length $n$ with autocorrelation function $\rho_k$ at lag $k$
has effective sample size:
\[
  \ess = \frac{n}{1 + 2\sum_{k=1}^{\infty} \rho_k}.
\]
For a chain with geometric autocorrelation decay $\rho_k \approx \rho^k$,
this simplifies to:
\[
  \ess \approx \frac{n(1 - \rho)}{1 + \rho}.
\]

The interpretation is direct: $\ess$ is the number of independent draws
that would provide the same estimator variance as the $n$-step chain.
An ensemble report should quote $\ess$, not $n$, when characterising
the precision of empirical quantile estimates.

\subsection{Redistricting-Specific ESS}

For \recom{} chains applied to redistricting, the autocorrelation
at lag 1 for Polsby-Popper is typically $\rho_1 \approx 0.85$--$0.95$,
reflecting the local-move nature of the proposal: merging two adjacent
districts and re-drawing the boundary changes at most a few percent of
tracts on each step.

With $\rho_1 = 0.90$ and geometric decay, a 10{,}000-step chain has:
\[
  \ess \approx \frac{10{,}000 \times 0.10}{1.90} \approx 526.
\]

This is a substantial reduction from the nominal 10{,}000 steps.
For Democratic seat counts, autocorrelation is typically lower
($\rho_1 \approx 0.70$--$0.80$) because seat counts change
discontinuously at district boundaries:
\[
  \ess(D\text{-seats}) \approx \frac{10{,}000 \times 0.25}{1.75} \approx 1{,}429.
\]

\subsection{State-Specific ESS Estimates}

Table~\ref{tab:ess-estimates} gives estimated $\ess$ for each study state
at 10{,}000 steps, based on empirical lag-1 autocorrelation.

\begin{table}[h]
\centering
\caption{Estimated ESS for 10{,}000 \recom{} steps. $\rho_1(\PP)$ from
         empirical Hamming autocorrelation; $\ess(\PP)$ computed via the
         geometric approximation $n(1-\rho_1)/(1+\rho_1)$.
         $\ess(D\text{-seats})$ uses a uniform assumed $\rho_1 = 0.75$
         for all states (D-seat autocorrelation is lower than PP because
         seat counts change discontinuously at district boundaries).}
\label{tab:ess-estimates}
\begin{tabular}{lcccc}
\toprule
State & $k$ & $\rho_1(\PP)$ & $\ess(\PP)$ & $\ess(D\text{-seats})$ \\
\midrule
NC & 14 & 0.87 & 695 & 1{,}429 \\
WI &  8 & 0.85 & 811 & 1{,}429 \\
GA & 14 & 0.88 & 638 & 1{,}429 \\
PA & 17 & 0.89 & 582 & 1{,}429 \\
TX & 38 & 0.92 & 417 & 1{,}429 \\
CA & 52 & 0.94 & 309 & 1{,}429 \\
\bottomrule
\end{tabular}
\end{table}

The pattern is clear: larger states have higher autocorrelation
(fewer tracts change on each step as a fraction of the total) and thus
lower $\ess$ for the same nominal step count.

\subsection{Minimum ESS for Reliable Quantile Estimates}

For a 95th-percentile tail estimate with standard error below 10 percentage
points of the outcome range ($\delta = 0.10$), the required $\ess$ is approximately:
\[
  \ess_{\min} \approx \frac{1}{4 \alpha (1-\alpha) \delta^2}
  = \frac{1}{4 \times 0.05 \times 0.95 \times (0.10)^2} \approx 526.
\]

We report the operational threshold as $\ess > 500$ because the diagnostic
tables are empirical and the threshold is meant as an evidentiary screen, not a
knife-edge mathematical identity. For claims about the 1st or 99th percentile,
the same formula with $\alpha=0.01$ gives
\[
  \ess_{\min,99} \approx
  \frac{1}{4 \times 0.01 \times 0.99 \times (0.10)^2} \approx 2{,}525.
\]
Those tail claims require much longer chains or a separate uncertainty
statement.

The 95th-percentile minimum is met at 10{,}000 steps for four of six study states:
NC (695), WI (811), GA (638), and PA (582) all exceed 526.
Texas (417) and California (309) fall below this threshold, requiring
approximately 15{,}000 and 20{,}000 steps respectively to achieve
10-percentage-point precision.
To meet a more stringent requirement of $\delta = 0.05$ (5 percentage
points), the required $\ess$ rises to approximately 2{,}105, requiring
roughly 60{,}000--90{,}000 steps for Texas and California.

This analysis motivates a statutory minimum of 10{,}000 steps for
four-state precision and 15{,}000--20{,}000 steps when Texas or California
are included in the ensemble, consistent with the MCMC practice in the
redistricting literature.
